\subsection{Ordinary real multiplication does not force a good canonical prime}

\paragraph{Finding A38 (failure of a universal canonical forcing statement).}
Let $C/\mathbf F_2$ be the ordinary genus-two curve of A13,
\[
                   y^2+y=x+x^{-1}+(x+1)^{-1}.
 \tag{BC.1}
\]
Its Jacobian has the integral Rosati-invariant $\mathcal O_5$
action proved in the supplement to A33. Among its $129$
$\mathbf F_{128}$-rational points, exactly three are Weierstrass.
For every one of the other $126$ points $R$,
\[
                  \operatorname{ord}([K_C-2R])=4159,
             \qquad \operatorname{ord}_{4159}(2)=2079
                       \quad\hbox{is odd}.
 \tag{BC.2}
\]
Consequently none of these canonical classes has a nonzero
$r$-primary component for any odd prime $r$ with even
$\operatorname{ord}_r(2)$. This remains true for each of these
fixed points after every constant extension. Thus integral
ordinary real multiplication by $\mathcal O_5$, even with a
non-Weierstrass prescribed point, does not force the sufficient
canonical-prime hypothesis of A29.

\begin{proof}
The three simple Artin--Schreier poles $0,1,\infty$ give
genus two and ordinarity; their ramification points are the
three rational Weierstrass points. The checked polynomial
in A13 is
\[
                         P_2(Z)=Z^4-Z^2+4.
 \tag{BC.3}
\]
Its roots occur in opposite pairs. Every odd power sum is
therefore zero, and the Frobenius point-count formula gives
$\#C(\mathbf F_{128})=128+1=129$
\cite[Theorem~11.1]{MilneJV}.
For completeness, let $\beta_1,\beta_2$ be the two roots of
$Z^2-Z+4$ and put $t_m=\beta_1^m+\beta_2^m$.
The recurrence $t_0=2$, $t_1=1$, $t_m=t_{m-1}-4t_{m-2}$
gives
\[
               (t_0,\ldots,t_7)=(2,1,-7,-11,17,61,-7,-251).
\]
Since seven is odd, pairing the opposite eigenvalues in
the Jacobian point-count product
\cite[Theorem~19.1(b)]{MilneAV} gives
\[
 \#J_C(\mathbf F_{128})
    =(1-\beta_1^7)(1-\beta_2^7)
    =4^7+1-t_7=16636=4\cdot4159.
 \tag{BC.4}
\]
Trial division through $\lfloor\sqrt{4159}\rfloor=64$
proves that $4159$ is prime. The exact residues
\[
 \begin{gathered}
 2079=3^3\cdot7\cdot11,\qquad
 2^{2079}\equiv1,\\
 2^{693}\equiv2554,\qquad
 2^{297}\equiv1245,\qquad
 2^{189}\equiv102\pmod{4159}
 \end{gathered}
 \tag{BC.5}
\]
certify the exact multiplicative order in (BC.2).
These are the arithmetic inputs already present in A13;
the retained additional certificate independently rechecks
them and directly enumerates all $126$ nonbranch points.

Write $G=J_C(\mathbf F_{128})$. The three distinct pairwise
differences of the rational Weierstrass points are all the
nonzero geometric points of $J_C[2]$, because the curve has
two-rank two. Thus $G[2]$ has order four and is
$(\mathbf Z/2)^2$. By (BC.4) it is already the entire
two-primary subgroup of $G$, which is therefore killed by two.
Choose a rational Weierstrass point $W$. Since $K_C\sim2W$,
\[
                         \kappa_R=[K_C-2R]=2[W-R].
 \tag{BC.6}
\]
Its two-primary part vanishes. If $R$ is non-Weierstrass,
this class is nonzero: otherwise $2R\sim K_C$ would be a
fibre of the canonical hyperelliptic pencil, forcing $R$ to
be a ramification point. Consequently its order is the
prime $4159$, proving (BC.2). Extension of constants preserves
the order of each fixed geometric divisor class, so it cannot
introduce another primary component.

An explicit example is available without a Jacobian-divisor
search. Let
\[
 k=\mathbf F_2[a]/(a^7+a+1),\qquad
 R=(a+1,\ a^5+a^4+a^2+a).
 \tag{BC.7}
\]
The finite certificate verifies the field, the curve equation,
and that this point is nonbranch; in binary coefficient
encoding its coordinates are $(3,54)$. Hence its canonical
class also has the exact order in (BC.2).
\end{proof}

\paragraph{What this does and does not decide after the quintic transfer.}
The example refutes the universal assertion that every ordinary
Rosati-$\mathcal O_5$ pointed genus-two curve has an A29
canonical detector. The example alone does not show that any of the
displayed points is the particular image of a prescribed split
point in a quintic transfer. A40 below supplies that inverse
dihedral construction with its pointing and field of definition.
A43 later supplies the stronger qualifications for the special
A13 inverse family: the sources have two-rank one and both
Weierstrass points have only bad canonical reflection images.
Those qualifications are proved there from the special Frobenius
identities, rather than inferred from real multiplication alone.

There is also no point-section counterexample here. A13 already
supplies a geometric theta witness of order $4159$ because
$\langle2,-1\rangle=(\mathbf Z/4159)^\times$, so A11
provides point separation after the indicated pointing extension.
The ordinary canonical criterion and that geometric theta
criterion have different hypotheses. Since (BC.1) is defined
over $\mathbf F_2$, its Frobenius twist is the same curve;
no identification over a nonperfect field is being assumed.
